\section{Funcionalidades}

A ferramenta tem seu foco em ataques a redes Wi-Fi. Mas, também conta com alguns scanners que são úteis para fornecer informações para serem usadas nos ataques. Elas são:


\subsection{\textbf{Ataque de Desautenticação}}
Esse ataque explora uma vulnerabilidade fundamental no protocolo 802.11, que gerencia as comunicações Wi-Fi. Ele é possível devido a duas características principais do protocolo:

\begin{itemize}
    \item Os quadros de gerenciamento, como os de autenticação e desautenticação, são enviados sem criptografia, exceto nas redes com WPA3 e WPA2 que usam o padrão 802.11w.

    \item O padrão Wi-Fi original não exige que os dispositivos verifiquem se um quadro de desautenticação foi enviado pelo ponto de acesso legítimo.
\end{itemize}

Um atacante pode enviar quadros de desautenticação forjados para um dispositivo (cliente) e/ou para o ponto de acesso (AP), fingindo ser a outra parte legítima da conexão. O dispositivo e/ou o AP aceita esses quadros falsificados como se fossem legítimos, pois não há um mecanismo de verificação de origem. Essa falha é usada em combinação com outros ataques como o ataque \texttt{kr00k} \cite{kr00k}, \texttt{KRACK} \cite{krack} e \texttt{evil-twin} \cite{evil-twin}.

O \offscan usa dessas falhas para executar o ataque. Ele cria os frames de desautenticação usando as informações passadas pela linha de comando e envia os frames em um loop infinito para o host alvo, se passando pelo AP, e para o AP, se passando pelo host alvo.

Essa falha foi descoberta em 2002 mas documentada em 2003. Sua primeira aparição foi no artigo \citet{80211-flaws} ainda em 2003. Mas essa falha ainda se mostra relevante pois são portados CVEs relacionados a ela, como o CVE-2026-26048 \cite{deauth-protecless}




\subsection{\textbf{Flooder de Beacons}}
Beacons são quadros (frames) enviados por roteadores de redes Wi-Fi para se apresentarem e enviarem informações necessárias para que dispositivos possam se conectar as suas redes sem fio.

O inundador de beacons cria vários desses beacons com informações falsas para que os dispositivos ao alcance recebem esses beacons e os processem como se fossem redes reais. Isso faz com que a lista de redes disponíveis fique maior e confunda os usuários, o que se enquadra em um ataque de confusão de SSID \cite{confusion-attack}. Uma das informações que o inundador usa é nome de rede (BSSID) que pode ser de uma rede real ou não. Para evitar filtragem de redes pela repetição do mesmo BSSID, o código cria variações do nome passado, alterando algumas letras de maiúscula para minúscula, e o inverso também. Assim como o ataque de desautenticação, esse flooder é comumente usado em conjunto com outras técnicas.

Esse flooder pode ser usado para ataques DoS em dispositivos IoT devido as limitações de hardware e/ou falhas no software. Esses ataques podem ocorrer por diferentes motivos como por consumo descontrolado de recursos \cite{cwe400}, alocação de recursos sem limites \cite{cwe770} e estouro de buffer baseado em pilha \cite{cwe121}.



\subsection{\textbf{Descoberta de dispositivos}}
Apesa de não ser um ataque, e sim uma técnica de reconhecimento \cite{ip-scanning}, ele é extremamente útil para serem usadados em ataque e/ou tomadas de decisão. Essa funcionalidade usa uma técnica de enviar pacotes especificos e analisar as respostas recebidas. No caso de uma escaneamento ativo.

A descoberta de dispositivos do \offscan envia pacotes ARP (para endereços locais), ICMP e TCP (por padrão, é possível escolher um ou todos os três serão usados para todos os IPs da rede, ou de um range especificado. O gerador recebe o CIDR (baseado na interface que será usada) ou um range específico e calcula todos os IPs possíveis e entrega um iterador para ser usado nos probes.

Durante o envio dos pacotes, um sniffer fica capturando as respostas em uma \textit{goroutine} (\textit{thread}) separada. Ele possui um filtro BPF (\textit{Berkley Packet Filter}) que busca por pacotes IPv4 e que contenham endereços de IP dentro do range especificado por um CIDR. Esse CIDR é criado basedo no primeiro e último IPs gerados pelo módulo \texttt{Ipv4Iter}. Assim, deixando a captura por pacotes mais especifica.



\subsection{\textbf{Descoberta de dispositivos na camada 2 (IEEE 802.11)}}

Essa técnica de descoberta de dispositivos tem objetivo semelhante ao da abordagem baseada em pacotes TCP/IP, porém difere na forma de operação, no alcance e no tipo de informação obtida.

Com a interface de rede em modo monitor, um sniffer captura dois tipos de quadros IEEE 802.11: \textit{beacons} e \textit{data frames}. A partir desses quadros, é possível descobrir quais dispositivos estão ativos no entorno e a quais redes eles estão conectados.

Os \textit{beacons} são enviados pelos pontos de acesso (APs) para anunciar suas redes. Deles, extraem-se apenas o SSID, o BSSID e o canal. Já os \textit{data frames} são trocados entre dispositivos e APs após o estabelecimento de uma conexão. Desses quadros, obtêm-se o endereço MAC do dispositivo e o BSSID.

Ao cruzar o BSSID presente nos \textit{data frames} com o BSSID dos \textit{beacons}, é possível identificar a qual rede cada AP pertence e quais dispositivos estão associados a ele.

Essa abordagem de descoberta é passiva e extremamente silenciosa. Ela só é viável porque os cabeçalhos dos quadros IEEE 802.11 não são criptografados — eles precisam ser legíveis para que os dispositivos possam verificar se cada quadro é destinado a eles ou deve ser descartado. Além disso, evitar a decriptação reduz o tempo de processamento e aumenta a vazão da rede.



\subsection{\textbf{Mapeador de redes Wi-Fi}}
O reconhecimento das redes Wi-Fi está cada vez mais sendo usado por atacantes por ser um vetor/porta inicial de ataque \cite{wifi-vec}. Buscar informações sobre as redes disponíveis pode revelar informações úteis para outras .

O mapedor de redes Wi-Fi tem como objetivo mostrar informações de cada rede wifi ao alcance. Essas informações são úteis para o ataque como o de desautenicação e outros.

As informações capturadas são:
\begin{itemize}
    \item BSSID (Endereço MAC)
    \item SSID (Nome da rede)
    \item Canal
    \item Frequência
    \item Padrão técnico (802.11X)
    \item Criptografia
\end{itemize}

Ele inicia o sniffer na interface e entra em um loop que fica alterando o canal da interface periodicamente. Após escanar cada interface, o sniffer é fechado e os pacotes são processados e as informações são formatadas e apresentadas.



\subsection{\textbf{Informações de Rede e Interface}}
Essa funcionalidade tem como objetivo mostrar algumas informações da interface e da rede, se estiver conectada, sem a necessidade de escanar a rede. As informações que são apresentadas são:

\begin{itemize}
    \item Nome da interface
    \item Estado da interface (UP, DOWN, ETC)
    \item Tipo (Ethernet, Wireless, etc)
    \item Endereço MAC
    \item Endereço de IP
    \item Endereço da rede e mascara (CIDR)
    \item Quantidade de IPs válidos
    \item MTU (Unidade Máxima de Transmissão)
    \item Endereço MAC do gateway
    \item Endereço de IP do gateway
\end{itemize}

Essas informações são buscadas de origens diferentes. Nome, estado, tipo, MAC e MTU são copiadas dos arquivos \texttt{/sys/class/net/<interface>}. Já os endereços de IP e MAC do gateway são copiados de arquivos que ficam nos diretórios \texttt{/proc/net/}. Já o IP e o endereço da rede com a máscara são conseguidos atravez de chamadas ao sistema usando a crate \texttt{libc}. E a quantidade de IPs válidos é calculada através do CIDR.

Essas informações são úteis para serem usadas pelos outros comandos e/ou para tomada de decisão.